
%(BEGIN_QUESTION)
% Copyright 2012, Tony R. Kuphaldt, released under the Creative Commons Attribution License (v 1.0)
% This means you may do almost anything with this work of mine, so long as you give me proper credit

100 ohm industrielle RTD-er er bygget av platinatråd. Kabler brukt for å koble platinabaserte RTD-elementer til fjernindikatorer eller transmittere er derimot nesten alltid kobberledere. Når kobberlederne kobles til RTD-en, dannes overganger mellom ulike metaller (kobber-platina).

\vskip 10pt

Til tross for disse overgangene mellom ulike metaller trenger vi aldri ``referansekrysskompensasjon'' i RTD-målesystemer, og vi trenger heller ikke spesielle kompensasjonskabler slik vi gjør i termoelementkretser. Forklar hvorfor metallovergangene i en RTD-krets ikke skaper noe måleproblem.

\underbar{file i00810}
%(END_QUESTION)





%(BEGIN_ANSWER)

De to kobber-platina-overgangene som oppstår når kobberledere kobles til platinatråden i RTD-sensoren, befinner seg på samme sted (ved RTD-sensoren) og har derfor samme temperatur. Polaritetene er motsatt rettet, så spenningsbidragene kansellerer hverandre.

\vskip 10pt

En annen forklaring er loven om mellomliggende metaller: når de to kobber-platina-overgangene har samme temperatur, blir platinatråden i RTD-elementet et mellomliggende metall mellom to kobberledere. Etter denne loven kan dette behandles som én kobber-kobber-overgang ved samme temperatur. En overgang mellom like metaller (kobber-kobber) gir ingen termoelektrisk spenning ved noen temperatur, og gir derfor ikke problemer i RTD-kretser slik referansekrysset gjør i termoelementkretser.

%(END_ANSWER)





%(BEGIN_NOTES)

{\bf Denne oppgaven er ment kun for eksamen, ikke for arbeidsark!}.

%(END_NOTES)
